\chapter{Введение: от архитектуры материи к динамике её рождения}

\section{Возвращение к исходной идее на новом уровне}

Том I исходил из общей идеи: частица может быть не добавленным к готовому
миру первичным объектом, а устойчивой модой единой структуры допустимых
состояний, переходов и ограничений. Он не предсказывал заранее конкретную
дефектную модель, а задал метод, по которому последующие модели должны
быть выведены, проверены и при необходимости отвергнуты. После
алгебраической, спектральной, топологической и вариационной работы Том V
получил строгую специализацию этой идеи: кандидат на материальную моду
может быть представлен дефектом неполного операторного замыкания.

В результате стало известно, как должен быть типизирован возможный
материальный дефект, какой стабильный класс он несёт и как измеряется его
минимальная неполнота замыкания. Но не было доказано, что пустой сектор
обязан породить такой дефект. Поэтому Том VI начинается не с очередной
классификации частицы и не с подбора её массы, а с более раннего вопроса:
может ли сама материальная ветвь возникнуть из нулевого состояния.

\section{Что Том VI наследует от первого тома}

Первый том является методологическим основанием программы, но не обязан
быть пророчеством её поздней математики. Поэтому дефектная конструкция,
класс пятнадцать и коэффициент $1/7$ не приписываются ему задним числом.
Они принадлежат результатам последующих томов.

От первого тома новый этап наследует пять требований:
\begin{enumerate}
 \item один источник должен заменять сумму независимых подгонок;
 \item фундаментальные данные должны быть отделены от производных
 интерпретаций;
 \item геометрическое, спектральное и корреляционное чтения должны
 оставаться совместимыми;
 \item рабочая гипотеза должна отличаться проверкой от нулевой;
 \item отсутствие различающего результата должно останавливать, а не
 маскировать программу.
\end{enumerate}

Таким образом, Том VI не переписывает исходную онтологию. Он проверяет,
оказался ли найденный в Томах IV--V дефект действительной реализацией
общей идеи устойчивой моды или только одним из математически допустимых
описаний.

\section{Граница, установленная Томом V}

Том VI принимает результаты предыдущего тома только с тем статусом,
который был фактически доказан. Исходный реестр разделяется на три слоя.

\subsection{Строго установленные входные данные}

\begin{enumerate}
 \item Переходы и их композиции допускают моритовский язык, а семейный и
 наблюдаемый углы помещаются в связывающую алгебру
 \begin{equation}
  \mathcal L(E)\simeq M_{35}(\mathbb C).
 \end{equation}

 \item Разность углов определяет нормированную суперразмерность
 \begin{equation}
  \frac{20-15}{35}=\frac17.
 \end{equation}

 \item Стабильный коэффициентный класс дефекта равен пятнадцати и имеет
 обменный вещественный подъём
 \begin{equation}
  \kappa_{15}=15\in KO_6(M_{105}(\mathbb C)_{\mathbb R}).
 \end{equation}

 \item Две ориентации образуют сопряжённую пару с числом намотки
 \begin{equation}
  (+15,-15),
 \end{equation}
 причём их обычный суммарный ориентированный индекс равен нулю.

 \item Необратимость представителя класса пятнадцать требует
 \begin{equation}
  \dim\ker T+\dim\operatorname{coker}T\geq15,
 \end{equation}
 а минимальный нормированный положительный дефицит равен
 \begin{equation}
  \delta_{\rm cl}=\frac{15}{105}=\frac17.
 \end{equation}

 \item Проекторный ёж со значениями в $\mathbb{RP}^2$ имеет два
 ориентированных спинорных подъёма. Они порождают хопфовы линии с
 $c_1=+1$ и $c_1=-1$, а после коэффициентного умножения --- локальные
 классы $+15$ и $-15$.

 \item Переход между нулевым и ненулевым топологическими секторами не
 может происходить внутри одной компоненты всюду обратимых операторов.
 Он обязан пересечь сингулярный слой, на котором возникает нулевая мода,
 теряется ранг или нарушается фредгольмова обратимость.

 \item Постоянный проекторный сектор существует и имеет нулевой заряд.
 Известный положительный петлевой функционал удовлетворяет
 \begin{equation}
  \mathcal S_{\rm loop}(I)=0,
  \qquad
  \min_{\operatorname{wind}=15}\mathcal S_{\rm loop}=\frac17.
 \end{equation}
 Следовательно, он минимизирует дефект только после фиксации сектора и не
 делает пустой вакуум неустойчивым.
\end{enumerate}

\subsection{Условно допустимые, но не выведенные элементы}

К началу Тома VI остаются математически допустимыми, но не доказанными:
единое пространство нулевого и дефектного секторов; родительский
функционал на этом пространстве; динамический закон прохождения через
необратимость; конечная энергия и радиус пары; физическая ориентация
детерминантной или пфаффиановой линии; космологическая история перехода;
отождествление классов $\pm15$ с конкретными частицами.

\subsection{Запрещённые подстановки}

В новом томе запрещается считать доказательством рождения материи:
само существование ненулевого $K$-класса; минимум действия внутри заранее
выбранного сектора; антропную условность; введённую вручную отрицательную
массу; выбранный без вывода член Скирма; фазу, тривиальную на полной
вещественной паре; или феноменологическое совпадение, найденное до
замыкания динамики.

\section{Центральная постановка задачи}

Пусть $I$ обозначает тождественный пустой сектор, а
$(V_{15},V_{-15})$ --- глобально нейтральную пару локально ненулевых
дефектов. Требуется определить, существует ли физически допустимый процесс
\begin{equation}
 I
 \longrightarrow
 T_{\rm sing}
 \longrightarrow
 (V_{15},V_{-15}),
 \label{eq:v6-zero-singular-pair}
\end{equation}
где $T_{\rm sing}$ является не формальной вставкой, а необходимым
необратимым слоем единой динамики. Потеря обратимости относится здесь к
конкретному операторному представителю и не должна автоматически
считаться разрушением инвариантного ядра общего источника. Если ядро не
сохраняется, процесс описывает переход к иной сущности, а не рождение
материи внутри одной теории.

\begin{definition}[Основная конфигурационная задача Тома VI]
Необходимо построить пространство конфигураций $\mathfrak C$, которое
одновременно содержит тождественный сектор, сингулярные переходные
операторы и конечные Real-пары дефектов, а также определить на нём один
родительский функционал
\begin{equation}
 \mathcal S_{\rm parent}\colon\mathfrak C\longrightarrow\mathbb R.
\end{equation}
Функционал не должен менять определение при переходе между секторами и не
должен содержать проектор на класс пятнадцать как внешнее условие.
\end{definition}

Главный локальный тест состоит в вычислении физического гессиана
тождественного сектора
\begin{equation}
 \Hess_I^{\rm phys}\mathcal S_{\rm parent},
\end{equation}
после удаления калибровочных, BRST/BV-точных, junk- и чисто координатных
направлений.

Отрицательная собственная мода
\begin{equation}
 \lambda_{\min}
 \bigl(\Hess_I^{\rm phys}\mathcal S_{\rm parent}\bigr)<0
 \label{eq:v6-negative-mode}
\end{equation}
является необходимым, но не достаточным признаком рождения материи. Нужно
дополнительно доказать, что нелинейное развитие этой моды проходит через
слой $T_{\rm sing}$ и насыщается конечной устойчивой парой, а не
неограниченным коллапсом, стенкой, калибровочным артефактом или иной
топологией.

\section{Рабочая и нулевая гипотезы}

\begin{description}
 \item[Рабочая гипотеза.] Тождественный сектор физически неустойчив, а
 его отрицательная мода нелинейно насыщается глобально нейтральной
 локальной парой $(V_{15},V_{-15})$, сохраняющей дефицит $1/7$.

 \item[Нулевая гипотеза.] После полного физического факторирования
 \begin{equation}
  \Hess_I^{\rm phys}\mathcal S_{\rm parent}\geq0,
 \end{equation}
 либо отрицательные направления не приводят к конечной паре требуемого
 класса. Тогда текущая архитектура не порождает материю самопроизвольно.
\end{description}

Нулевая гипотеза имеет равный исследовательский статус с рабочей. Том VI
не обязан подтвердить рождение материи; он обязан провести различающий
тест.

\section{Основные задачи Тома VI}

\begin{enumerate}
 \item Определить минимальное общее пространство операторов, содержащее
 $I$, $T_{\rm sing}$ и $(V_{15},V_{-15})$.
 \item Типизировать допустимые пути потери обратимости и связать их со
 спектральным потоком, ядром и коядром переходного оператора.
 \item Вывести один родительский функционал без ручной фиксации
 топологического сектора.
 \item Построить физический фактор пространства вариаций и вычислить
 полный гессиан тождественного вакуума.
 \item Если отрицательная мода существует, получить её нелинейное
 продолжение и проверить, является ли конечный исход Real-парой
 $+15/-15$.
 \item Доказать или опровергнуть конечность энергии, локализацию и
 существование устойчивого радиуса пары.
 \item Проверить устойчивость относительно всех допустимых возмущений,
 а не только радиального анзаца.
 \item Построить квантовую меру, проверить determinant/Pfaffian line,
 аномалии и правило суммирования по ориентациям.
 \item Лишь после динамического замыкания определить допустимую карту к
 массам, размерам, взаимодействиям и космологической плотности дефектов.
 \item Сформулировать слепые наблюдательные тесты до сравнения с данными.
\end{enumerate}

\section{Критерий успеха}

Основная задача считается положительно закрытой только при одновременном
выполнении следующих условий:
\begin{enumerate}
 \item нулевой сектор и дефектная пара принадлежат одной определённой
 конфигурационной задаче;
 \item родительский функционал и мера одинаково определены на обоих
 секторах и переходном слое;
 \item после физического факторирования существует настоящая отрицательная
 мода нулевого сектора;
 \item её нелинейное развитие проходит через необходимую потерю
 обратимости;
 \item конечный исход является локализованной Real-парой с классами
 $+15/-15$, нулевым полным ориентированным зарядом и положительным
 дефицитом $1/7$;
 \item энергия ограничена снизу, а размер пары конечен и определяется
 функционалом;
 \item результат устойчив к замене координат, калибровки, регуляризации
 и допустимого представителя стабильного класса.
\end{enumerate}

\section{Критерии провала и остановки}

\begin{nogocriterion}[Стабильный пустой вакуум]
Если полный физический гессиан неотрицателен, текущий родитель не порождает
материю. Переход к вычислению масс и подбору феноменологии запрещается.
\end{nogocriterion}

\begin{nogocriterion}[Ложная отрицательная мода]
Если отрицательное направление является калибровочным, junk-, BRST/BV-
точным, координатным или исчезает после полного факторирования, оно не
считается физической неустойчивостью.
\end{nogocriterion}

\begin{nogocriterion}[Неправильное нелинейное насыщение]
Если отрицательная мода ведёт к неограниченному действию, делокализации,
стенке, одиночному нескомпенсированному заряду или сектору, не связанному
с $\kappa_{15}$, механизм рождения материи не закрыт.
\end{nogocriterion}

\begin{nogocriterion}[Ручной выбор материи]
Если класс пятнадцать, пространственный масштаб или стабилизирующий член
вводятся только внешним ограничением, результат является условной моделью
дефекта, а не выводом рождения материальной ветви.
\end{nogocriterion}

\section{Порядок исследования}

Работа должна идти в следующем порядке:
\begin{equation}
 \boxed{\begin{gathered}
 \text{пространство конфигураций}
 \to \text{общий функционал}
 \to \text{физический гессиан},\\
 \text{нелинейная пара}
 \to \text{квантовая мера}
 \to \text{физическое чтение}
 \end{gathered}}
\end{equation}

До закрытия первых четырёх звеньев запрещается объявлять классы
$\pm15$ частицами Стандартной модели, трактовать $1/7$ как массу или
вероятность и использовать наблюдательные числа для настройки
родительского функционала.

\section{Итоговая формулировка программы}

Том V ответил на вопрос, какой строгий операторно-топологический след
оставляет материальный дефект, если ненулевой сектор уже существует.
Том VI должен ответить на более глубокий вопрос: существует ли единый
закон, при котором пустой сектор сам теряет устойчивость и рождает
глобально нейтральную пару таких дефектов.

Положительный ответ откроет переход от архитектуры к физике рождения
материи. Отрицательный ответ будет не менее содержателен: он докажет, что
нынешняя родительская конструкция описывает допустимые дефекты, но не
объясняет, почему мир выбирает материальную ветвь.